\begin{table}[t]
\centering
\caption{Standardized Effect Sizes}
\label{tab:sde}
\scriptsize\begin{tabular}{lcccccc}
\toprule
Outcome & $\hat{\beta}$ & SE & SD($Y$) & SDE & SE(SDE) & Classification \\
\midrule
\emph{Panel A: Pooled (strongest specification)} \\
\addlinespace
Total employment (\% $\Delta$) & 131.25 & 37.24 & 5.00 & 0.388 & 0.110 & Large \\
Unemployment rate ($\Delta$ pp) & 34.08 & 22.43 & 3.84 & 0.131 & 0.086 & Moderate \\
\addlinespace
\emph{Panel B: Heterogeneous (Eurozone vs.~non-Eurozone)} \\
\addlinespace
Total empl., Eurozone & 110.76 & 37.83 & 5.00 & 0.327 & 0.112 & Large \\
\bottomrule
\end{tabular}
\par\vspace{0.3em}
{\scriptsize \textit{Notes:}
\textbf{Country:} European Union (27 member states).
\textbf{Research question:} Does exposure to CRD~IV/Basel~III bank branch closures cause regional employment decline in EU NUTS2 regions?
\textbf{Policy mechanism:} CRD~IV (2013) and CRR imposed capital adequacy, liquidity coverage, and net stable funding requirements on EU banks, disproportionately burdening small cooperative and retail banks, triggering closure of over 110,000 bank offices (2008--2024).
\textbf{Outcome definition:} Total employment change as percentage deviation from 2008 baseline, from Eurostat LFS regional tables, ages 15--64.
\textbf{Treatment:} Continuous; pre-shock (2008) NACE Rev.~2 Section~K employment share interacted with post-CRD~IV indicator (2014+).
\textbf{Data:} Eurostat NUTS2 panels, 2008--2023; 219 regions, 27 countries; strongest spec.\ excludes 16 capital regions ($N = 3{,}038$).
\textbf{Method:} OLS with NUTS2 and country-by-year FE; SEs clustered at country level (27 clusters).
\textbf{Sample:} NUTS2 regions with non-missing 2008 NACE~K employment; capitals excluded to mitigate composition bias.
SDE $= \hat{\beta} \times \text{SD}(X) / \text{SD}(Y)$ where SD($Y$) is pre-treatment standard deviation.
Classification refers to magnitude, not statistical significance: Large ($|$SDE$| > 0.15$), Moderate ($0.05$--$0.15$), Small ($0.005$--$0.05$), Null ($< 0.005$).}
\end{table}
